\documentclass[a4paper,14pt,oneside]{extarticle}
\usepackage[left=30mm,right=15mm,top=20mm,bottom=20mm]{geometry}
\usepackage{fontspec}
\usepackage[russian]{babel}
\usepackage{setspace}
\usepackage{titlesec}
\usepackage{indentfirst}

\IfFontExistsTF{Times New Roman}
  {\setmainfont{Times New Roman}}
  {\setmainfont{Liberation Serif}}
\onehalfspacing
\setlength{\parindent}{12.5mm}
\setlength{\parskip}{0pt}

\titleformat{\section}[block]{\centering\fontsize{16bp}{19bp}\selectfont\bfseries}{}{0pt}{}
\titlespacing*{\section}{0pt}{0pt}{\baselineskip}

\begin{document}
\fontsize{14bp}{17bp}\selectfont
\thispagestyle{empty}
\section*{Паспорт НИР}

\textbf{Тема.} Сравнение графовых нейронных сетей для предсказания недостающих связей в графе научных цитирований.

\textbf{Актуальность.} Пропуски в цитированиях затрудняют поиск связанных работ.
  Модель может предлагать ссылки для последующей ручной проверки.

\textbf{Объект исследования.} Ориентированный граф, вершины которого представляют научные статьи,
  а рёбра~--- цитирования.

\textbf{Предмет исследования.} Влияние агрегации признаков соседних вершин на качество предсказания связей.

\textbf{Цель.} Проверить, лучше ли GraphSAGE ранжирует ссылки по сравнению
  с эвристикой на основе степеней вершин и матричной факторизацией.

\textbf{Задачи.} Построить синтетический граф, разделить данные по времени,
  обучить модели и сравнить качество и время работы.

\textbf{Методика.} Для обучения использовать цитирования до выбранной даты.
  При проверке ранжировать более поздние цитирования среди статей, доступных на момент прогноза.
  Метрики~--- средний обратный ранг (MRR) и доля попаданий в первые десять результатов (Hits@10).

\textbf{Ожидаемый результат.} Воспроизводимая методика оценки, прототип и сравнение методов.

\textbf{Практическое применение.} Ранжирование возможных ссылок в системах поиска научных публикаций.

\textbf{Статус.} Учебный пример исследования на синтетических данных.
\end{document}
